% Generated draft: requires mathematical and source review.
\documentclass[11pt]{article}
\usepackage[margin=25mm]{geometry}
\usepackage{fontspec}
\setmainfont{Latin Modern Roman}
\usepackage{amsmath,amssymb,mathtools,unicode-math}
\setmathfont{Latin Modern Math}
\usepackage{enumitem,needspace,xurl,hyperref}
\setlength{\parindent}{0pt}
\setlength{\parskip}{5pt}
\setlength{\emergencystretch}{3em}
\newcommand{\N}{\mathbb{N}}
\newcommand{\Z}{\mathbb{Z}}
\newcommand{\Q}{\mathbb{Q}}
\newcommand{\R}{\mathbb{R}}
\newcommand{\C}{\mathbb{C}}
\newcommand{\F}{\mathbb{F}}
\newcommand{\A}{\mathbb{A}}
\newcommand{\PP}{\mathbb{P}}
\newcommand{\E}{\mathbb{E}}
\newcommand{\eps}{\varepsilon}
\newcommand{\dd}{\,\mathrm{d}}
\begin{document}
\section*{UnsolvedMath ID: 30004734}
\textbf{Source problem code:} catalogue entry
\subsection{Definitions and mathematical statement}
A periodic billiard trajectory is a trajectory consisting of straight segments in a billiard domain, obeying the specular reflection law at each regular boundary collision, and returning to its initial position and direction after finitely many reflections. Its length is the sum of the lengths of its segments over the specified period. For a billiard domain \(\Omega\) and \(L>0\), write \(\mathcal P_L(\Omega)\) for the collection of periodic billiard trajectories of length \(L\). The source does not specify whether trajectories are identified under changes of starting point or reversal, or whether iterates are included.

To describe the spectral terminology conditionally, suppose that \(\Omega\subset\R^2\) is a bounded domain with sufficiently regular boundary, and that a self-adjoint boundary condition \(B\) has been chosen for the nonnegative Laplacian \(\Delta_{\Omega,B}\). If its eigenvalues, counted with multiplicity, are \(\lambda_j\), two usual trace conventions are the distributional traces
\[
U_{\Omega,B}(t)=\operatorname{Tr}\bigl(e^{it\sqrt{\Delta_{\Omega,B}}}\bigr)=\sum_j e^{it\sqrt{\lambda_j}},
\qquad
C_{\Omega,B}(t)=\operatorname{Tr}\bigl(\cos(t\sqrt{\Delta_{\Omega,B}})\bigr)=\sum_j\cos(t\sqrt{\lambda_j}).
\]
A wave-trace expansion at a length \(L\) describes the singular asymptotic behavior of the chosen trace distribution near \(t=L\). These definitions provide a possible conventional setting; the source does not specify the domain class, boundary condition, trace convention, or type of asymptotic expansion.

\textbf{Source question.} Can a set of periodic billiard orbits, for example a positive-measure Cantor set of a fixed length, produce the wave-trace expansion of an ellipse?

In this terminology, the proposed example concerns a Cantor set \(K\subseteq\mathcal P_L(\Omega)\) at a common length \(L\), with positive measure in an unspecified ambient space and measure. The intended question is whether such a collection can produce the wave-trace expansion associated with an ellipse. The source does not identify the comparison ellipse, the relevant length or lengths, the meaning of the contribution of an individual collection of orbits, or the required extent of agreement of expansions. It also does not explicitly require \(\Omega\) to be nonelliptical. Consequently, a unique precise quantified existence or impossibility statement cannot be extracted without further hypotheses; these qualifications remain unresolved in this statement-only draft.
\subsection{Short English statement}
The question asks whether a collection of periodic billiard trajectories, possibly a positive-measure Cantor set whose trajectories all have the same length, can account for the wave-trace expansion of an ellipse. The measure, billiard class, boundary conditions, and meaning of matching the expansion are unspecified, so the intended precise formulation requires review.
\subsection{Sources}
\begin{itemize}
\item Local UnsolvedMath registry, supplied complete source record: permanent ID \(30004734\), source problem code \texttt{catalogue entry}, title \emph{Periodic Billiards and the Wave Trace of Ellipses}; collection: Oberwolfach Reports Open Problems. No individual proposer is identified. The supplied proposed year is 2021.
\item Supplied report citation: \emph{Dynamische Systeme (2022)}, Oberwolfach Reports, DOI \texttt{10.4171/owr/2021/33}, pp.\ 1735--1803. \url{https://doi.org/10.4171/owr/2021/33}. The background describes these reports as documenting mathematical workshops and their open-problem discussions.
\item Supplied literature source, labelled ``analytic symmetric rigidity'': \url{https://annals.math.princeton.edu/2009/170-1/p06}.
\item Supplied literature source, labelled ``local uniqueness of almost every ellipse'': \url{https://doi.org/10.1007/s00222-025-01397-y}.
\item Supplied literature source, labelled ``formal ellipse rigidity'': \url{https://msp.org/apde/2012/5-5/apde-v5-n5-p06-s.pdf}.
\item Supplied literature assessment, checked at \texttt{2026-08-22}: the record labels the problem partially solved and states that spectral rigidity is proved in analytic symmetry classes, that 2025 work gives local uniqueness for almost every ellipse, and that full rigidity of every ellipse without symmetry assumptions remains open. This is supplied dated curation derived from the OWR curation, not an independently verified assessment. The source explicitly requests independent verification of time-sensitive status and cited claims.
\end{itemize}
\end{document}
